\documentclass[aspectratio=169]{beamer}
\usepackage[utf8]{inputenc}
\usepackage[T1]{fontenc}
\usepackage[brazilian]{babel}
\usepackage{lmodern}
\usepackage{xcolor}
\usetheme{Madrid}
\setbeamertemplate{navigation symbols}{}
\setbeamercolor{structure}{fg=blue!55!black}
\setbeamercolor{block title}{bg=blue!12,fg=black}
\setbeamercolor{block body}{bg=black!3,fg=black}

\title{Conhecimento e expressão na sociedade maia}
\subtitle{Relações históricas sem equivalências anacrônicas}
\author{Carlos E. L. Pereira e David F. Oliveira}
\institute{IMPA Tech}

\begin{document}
\begin{frame}[plain]
  \titlepage
\end{frame}

\begin{frame}{Uma pergunta, três objetos}
\begin{block}{Questão central}
Como observação, representação e vida social se relacionavam no mundo maia, sem presumir que as categorias modernas de ciência, arte e religião fossem fronteiras universais?
\end{block}

\vspace{0.2cm}
\begin{columns}[T,onlytextwidth]
  \column{0.32\textwidth}
  \textbf{Castillo de Kukulcán}\\[0.08cm]
  Luz e sombra observáveis; intenção equinocial específica não demonstrada.

  \column{0.32\textwidth}
  \textbf{Códice de Dresden}\\[0.08cm]
  Registros calendáricos e astronômicos; alcance preditivo local não verificado no corpus.

  \column{0.32\textwidth}
  \textbf{Bonampak}\\[0.08cm]
  Imagem e escrita em narrativas sociais e políticas; leitura astronômica permanece hipótese.
\end{columns}

\vfill
\begin{alertblock}{Síntese}
Práticas hoje classificadas separadamente podiam estar material e institucionalmente imbricadas. Isso não significa que fossem idênticas nem que os maias compartilhassem nossas categorias.
\end{alertblock}
\end{frame}
\end{document}
